%==============================================================
% ABSTRAK (Bahasa Indonesia)
%==============================================================

\chapter*{ABSTRAK}

\noindent
ALDI PRAMUDYA. Deteksi Perubahan Deforestasi Menggunakan \textit{Siamese U-Net} Pada Citra Satelit Multispektral.
Dibimbing oleh AHMAD RIDHA dan MEDRIA KUSUMA DEWI HARDIHENATA.

Deforestasi merupakan ancaman serius bagi keanekaragaman hayati dan keseimbangan iklim global, terutama di hutan tropis Indonesia. Metode deteksi perubahan tutupan lahan konvensional seringkali kurang efektif dalam menangani kompleksitas citra multispektral dan variabilitas kondisi atmosfer. Penelitian ini bertujuan mengembangkan model deteksi perubahan tutupan lahan akibat deforestasi secara otomatis menggunakan arsitektur \textit{Siamese U-Net} berbasis \textit{deep learning} pada citra Sentinel-2. Model dirancang dengan dua cabang \textit{encoder} berbobot bersama (\textit{weight sharing}) untuk mengekstraksi fitur dari sepasang citra multi-temporal, lalu digabungkan pada modul \textit{decoder} guna menghasilkan peta perubahan biner. Performa model akan dievaluasi menggunakan metrik \textit{F1-score}, \textit{Intersection over Union} (IoU), dan akurasi keseluruhan, serta dibandingkan dengan beberapa metode pembanding. Penelitian ini diharapkan dapat menunjukkan keunggulan pendekatan \textit{Siamese U-Net} dalam mendeteksi deforestasi secara akurat guna mendukung sistem pemantauan hutan berbasis \textit{remote sensing} secara operasional.

\vspace{0.3cm}
\noindent
Kata kunci: deforestasi, deteksi perubahan, \textit{deep learning}, Sentinel-2, \textit{Siamese U-Net}
%==============================================================
% ABSTRACT (Bahasa Inggris)
%==============================================================

\vspace{1cm}
\begingroup
\let\clearpage\relax
\let\cleardoublepage\relax
\chapter*{ABSTRACT}
\endgroup


\noindent
ALDI PRAMUDYA. Deforestation Change Detection Using Siamese U-Net On Multispectral Satellite Imagery.
Supervised by AHMAD RIDHA and MEDRIA KUSUMA DEWI HARDIHENATA.

Deforestation poses a serious threat to global biodiversity and climate stability, particularly in Indonesian tropical forests. Conventional land cover change detection methods are often inadequate in handling the complexity of multispectral imagery and atmospheric variability. This study aims to develop an automated land cover change detection model for deforestation using a Siamese U-Net deep learning architecture on Sentinel-2 satellite imagery. The model is designed with two weight-sharing encoder branches to extract features from a pair of multi-temporal images, which are then fused within a decoder module to generate a binary change map. Model performance will be evaluated using F1-score, Intersection over Union (IoU), and overall accuracy, and compared against alternative methods. The results of this study are expected to demonstrate the superiority of the Siamese U-Net approach in accurately detecting deforestation to support operational remote sensing-based forest monitoring systems.

\vspace{0.3cm}
\noindent
\textit{Keywords}: change detection, deep learning, deforestation, Sentinel-2, Siamese U-Net
